\documentclass[journal,twoside,web]{ieeecolor}

\usepackage{generic}
\usepackage{cite}
\usepackage{amsmath,amssymb,amsfonts}
\usepackage{algorithmic}
\usepackage{graphicx}
\usepackage{algorithm,algorithmic}
\usepackage{hyperref}
\hypersetup{hidelinks=true}
\usepackage{textcomp}
\usepackage{booktabs}
\usepackage{multirow}
\usepackage{array}
\usepackage{enumitem}

\def\BibTeX{{\rm B\kern-.05em{\sc i\kern-.025em b}\kern-.08em
T\kern-.1667em\lower.7ex\hbox{E}\kern-.125emX}}

\markboth{\hskip25pc IEEE TRANSACTIONS ON MEDICAL IMAGING}
{Author \MakeLowercase{\textit{et al.}}: Sharing or Specializing?}

\begin{document}

\title{Sharing or Specializing? A Systematic Study of
Parameter-Efficient Vision--Language Model Adaptation for
Heterogeneous Medical Image Analysis}

\author{
First A. Author,
Second B. Author,
Third C. Author,
et al.
\thanks{Manuscript received Month XX, 2026; revised Month XX, 2026.}
\thanks{Author affiliations and funding information will be added before
submission.}
}

\maketitle


% ============================================================
\begin{abstract}

General-purpose vision--language models (VLMs) provide a unified interface
for heterogeneous medical image analysis, but adapting a single pretrained
model to tasks with fundamentally different visual, semantic, and output
requirements raises an unresolved design question: when should adaptation
parameters be shared across tasks, and when should they be specialized?
This study systematically investigates parameter-efficient adaptation across
classification, multi-label classification, detection, counting, regression,
and report generation spanning diverse medical imaging domains and
modalities. We compare fully shared LoRA, task-specific LoRA, partially
shared HydraLoRA, soft adapter composition, modality-aware adaptation,
latent LoRA mixture-of-experts, spectral expert methods, and alternative
specialization locations under controlled experimental settings. We further
examine cross-task transfer, supervision scaling, routing behavior,
representation similarity, domain-level specialization, and interactions
with practical design choices including image resolution, adapter rank,
preprocessing, augmentation, prompting, multi-image handling, decoding,
and backbone choice. The objective is to identify the factors associated
with successful parameter sharing and specialization when general-purpose
VLMs are adapted to heterogeneous medical imaging tasks. Final conclusions
and design recommendations will be based only on experimentally verified
results.

\end{abstract}

\begin{IEEEkeywords}
Medical image analysis, mixture of experts, parameter-efficient fine-tuning,
task specialization, transfer learning, vision--language models.
\end{IEEEkeywords}


% ============================================================
\section{Introduction}
\label{sec:introduction}

\IEEEPARstart{G}{eneral-purpose} vision--language models (VLMs) have emerged
as a promising unified framework for medical image analysis. A single model
can potentially support disease classification, multi-label recognition,
lesion localization, quantitative estimation, counting, and clinical report
generation through a common image--text interface. This flexibility is
attractive for developing general-purpose medical artificial intelligence
systems, but it also introduces substantial heterogeneity in visual
requirements, output spaces, and supervision regimes.

Parameter-efficient fine-tuning (PEFT), particularly Low-Rank Adaptation
(LoRA), provides an efficient mechanism for adapting large pretrained VLMs
without updating the full model. However, heterogeneous medical tasks create
an important design choice. At one extreme, a single shared adapter can be
trained across all tasks, maximizing data sharing but potentially increasing
cross-task interference. At the other extreme, independent task-specific
adapters can reduce interference but restrict each adapter to task-specific
supervision. Between these extremes lie partially shared, softly composed,
modality-aware, and mixture-of-experts formulations.

Existing approaches therefore make different assumptions about how
specialization should be structured. Some explicitly separate downstream
tasks, others share low-rank components, some introduce modality-specific
experts, and others rely on learned routing to discover latent expert
structure. However, these approaches are generally evaluated independently,
making it difficult to determine which properties of a medical task favor
sharing and which favor specialization.

The central question of this study is therefore:

\begin{quote}
\emph{When adapting a general-purpose VLM to heterogeneous medical imaging
tasks, which adaptation parameters should be shared and which should be
specialized?}
\end{quote}

Rather than assuming that any one specialization mechanism is optimal, this
work evaluates a spectrum of adaptation strategies under controlled
experimental conditions. We investigate the roles of task identity,
task-specific data availability, output structure, imaging modality, medical
domain, task compatibility, representation similarity, and specialization
location.

Our experimental setting contains six heterogeneous medical task families:
classification, multi-label classification, detection, counting, regression,
and report generation. These tasks provide a useful testbed because they
differ substantially in output format and visual reasoning requirements.

The study is organized around the following research questions:

\begin{enumerate}[label=\textbf{RQ\arabic*.}]
    \item Does explicit task specialization outperform shared adaptation?

    \item How does task-specific supervision affect the
    sharing--specialization trade-off?

    \item Can partially shared or dynamically composed adapters improve over
    fully shared and fully specialized adaptation?

    \item What should experts specialize on: task, modality, medical domain,
    or latent representation?

    \item Can learned routers discover meaningful specialization?

    \item What predicts positive or negative transfer between medical tasks?

    \item Where should specialization be introduced within the VLM?

    \item How do practical training and inference choices interact with
    specialization?
\end{enumerate}


% ============================================================
\section{Related Work}
\label{sec:related}

\subsection{General and Medical Vision--Language Models}

This section will review general-purpose VLMs and medically adapted or
medically pretrained VLMs, with emphasis on their use for heterogeneous
medical image analysis and medical visual question answering.

A key distinction will be made between general-purpose VLMs that are
subsequently adapted to medical tasks and models whose pretraining or
instruction-tuning data are specifically medical.

\subsection{Parameter-Efficient Adaptation}

This section will review LoRA, QLoRA, adapter tuning, task-specific
adapters, partially shared low-rank adaptation, and related PEFT approaches.

Particular attention will be given to how existing methods distribute
trainable parameters between shared and specialized components.

\subsection{Multitask Learning and Negative Transfer}

This section will review parameter sharing, task interference, gradient
conflict, positive transfer, and negative transfer in multitask learning.

These concepts will be related to the heterogeneous output structures
encountered in medical VLM adaptation.

\subsection{Mixture-of-Experts and Routed LoRA}

This section will review mixture-of-experts architectures, LoRA-MoE,
sparse expert routing, spectral expert methods, HydraLoRA, and
medical-domain expert models.

\subsection{Cross-Task and Cross-Modality Generalization}

This section will review prior studies evaluating transfer and
generalization across medical imaging modalities and clinical task types.


% ============================================================
\section{Problem Formulation}
\label{sec:problem}

Let a pretrained VLM be parameterized by $\theta_0$. For task
$t\in\mathcal{T}$, parameter-efficient adaptation introduces an update
$\Delta\theta_t$ such that

\begin{equation}
f_t(x,q)=f(x,q;\theta_0+\Delta\theta_t),
\end{equation}

where $x$ denotes the medical image and $q$ denotes the textual instruction.

For a LoRA-adapted linear transformation,

\begin{equation}
W_t=W_0+\Delta W_t,
\end{equation}

with

\begin{equation}
\Delta W_t=B_tA_t,
\end{equation}

where $A_t\in\mathbb{R}^{r\times d}$ and
$B_t\in\mathbb{R}^{d'\times r}$.

The central problem is to determine how the adaptation parameters
$\Delta W_t$ should be structured and shared across heterogeneous tasks.


% ============================================================
\subsection{The Sharing--Specialization Spectrum}

We organize the evaluated adaptation approaches according to the degree and
form of parameter sharing.

\subsubsection{Shared LoRA}

A single adapter is trained across all task families:

\begin{equation}
\Delta W_t=\Delta W_{\mathrm{shared}},
\qquad \forall t.
\end{equation}

\subsubsection{Task-Specific LoRA}

Independent low-rank parameters are trained for each task:

\begin{equation}
\Delta W_t=B_tA_t.
\end{equation}

\subsubsection{HydraLoRA}

A shared low-rank component is combined with task-specific components:

\begin{equation}
\Delta W_t=B_tA_s.
\end{equation}

\subsubsection{Soft Adapter Composition}

Multiple adapters are combined through learned or predefined weights:

\begin{equation}
\Delta W(x,q)=
\sum_{k=1}^{K}
\pi_k(x,q)\Delta W_k.
\end{equation}

\subsubsection{Modality-Aware Adaptation}

Task and modality-specific adaptation can be combined as

\begin{equation}
\Delta W(x,q)
=
\Delta W_t+\lambda_m\Delta W_m,
\end{equation}

where $\Delta W_m$ represents a modality-specific update.

\subsubsection{Latent LoRA Mixture-of-Experts}

A learned router selects or combines latent experts:

\begin{equation}
\Delta W(x,q)=
\sum_{k\in\mathrm{TopK}(g(h))}
\pi_k(h)B_kA_k,
\end{equation}

where $h$ denotes the routing representation and $g(\cdot)$ denotes the
router.

\subsubsection{Spectral Expert Adaptation}

Spectral expert approaches construct or initialize experts from distinct
spectral components of pretrained weights and combine them through learned
sparse routing.


% ============================================================
\section{Experimental Framework}
\label{sec:framework}

\subsection{Medical Task Families}

We evaluate six heterogeneous medical image-analysis task families.

\begin{table}[!t]
\caption{Task heterogeneity considered in the study.}
\label{tab:tasks}
\centering
\setlength{\tabcolsep}{3pt}
\begin{tabular}{lll}
\toprule
Task & Output Structure & Main Capability \\
\midrule
Classification & Class label & Recognition \\
Multi-label & Label set & Recognition/calibration \\
Detection & Bounding boxes & Localization \\
Counting & Integer & Quantification \\
Regression & Continuous scalar & Estimation \\
Report generation & Free text & Clinical generation \\
\bottomrule
\end{tabular}
\end{table}


% ============================================================
\subsection{Datasets and Medical Domains}

This subsection will describe:

\begin{itemize}
    \item constituent medical datasets,
    \item imaging modalities,
    \item medical domains,
    \item number of images,
    \item number of question--answer pairs,
    \item task distribution,
    \item class distributions where applicable.
\end{itemize}

The benchmark will be presented as a heterogeneous medical imaging testbed
rather than primarily as a challenge dataset.


% ============================================================
\subsection{Primary Backbone}

The final Qwen VLM will serve as the primary experimental platform.

The following implementation details will be reported:

\begin{itemize}
    \item exact backbone version,
    \item quantization configuration,
    \item frozen components,
    \item trainable components,
    \item LoRA target layers,
    \item LoRA rank,
    \item LoRA scaling,
    \item dropout,
    \item optimizer,
    \item learning rate,
    \item number of epochs,
    \item image resolution,
    \item inference configuration.
\end{itemize}


% ============================================================
\subsection{Secondary Backbone}

One additional VLM will be used to assess whether the principal observations
depend on the primary backbone.

The secondary evaluation will focus only on:

\begin{enumerate}
    \item Shared LoRA versus TaskLoRA,
    \item TaskLoRA versus HydraLoRA,
    \item supervision scaling.
\end{enumerate}

A medically pretrained VLM will be considered if computationally and
architecturally feasible.


% ============================================================
\subsection{Evaluation Protocol}

Controlled experiments will use fixed image-disjoint
A$\rightarrow$B and B$\rightarrow$A splits.

Results from organizer-provided validation or test sets, where applicable,
will be reported separately from internal model-selection experiments.

Historical experiments that were performed using different optimization
budgets will not be used as direct architectural comparisons unless rerun
under matched conditions.


% ============================================================
\subsection{Evaluation Metrics}

Task-specific metrics will be used:

\begin{itemize}
    \item Classification: balanced accuracy;
    \item Multi-label classification: F1;
    \item Detection: F1 at the specified intersection-over-union threshold;
    \item Counting: mean absolute error;
    \item Regression: mean absolute error;
    \item Report generation: official report-generation metric.
\end{itemize}

For principal comparisons, 95\% bootstrap confidence intervals will be
reported.


% ============================================================
\section{Study 1: Shared Versus Task-Specific Adaptation}
\label{sec:study1}

The first study will establish the two endpoints of the
sharing--specialization spectrum.

We will compare:

\begin{enumerate}
    \item the base VLM,
    \item a single shared LoRA adapter,
    \item independent task-specific LoRA adapters.
\end{enumerate}

The analysis will examine:

\begin{itemize}
    \item per-task performance,
    \item aggregate performance,
    \item cross-task interference,
    \item output-format errors,
    \item variation across medical domains.
\end{itemize}

The purpose of this study is to determine which task families benefit from
explicit specialization and which can be effectively represented through a
shared adaptation.


% ============================================================
\section{Study 2: Cross-Task Transfer}
\label{sec:study2}

To characterize task compatibility directly, each task-specific adapter
trained on source task $i$ will be evaluated on target task $j$.

We define the transfer score as

\begin{equation}
C_{ij}
=
M_j(\Delta W_i)
-
M_j(\Delta W_{\mathrm{base}}),
\end{equation}

where $M_j$ denotes the target-task metric after orienting metrics such that
larger values represent better performance.

A complete source--target task transfer matrix will be constructed.

Transfer patterns will be analyzed according to:

\begin{itemize}
    \item output-structure similarity,
    \item task family,
    \item source-task data size,
    \item target-task data size,
    \item modality overlap,
    \item medical-domain overlap.
\end{itemize}


% ============================================================
\section{Study 3: Data Availability and Parameter Sharing}
\label{sec:study3}

A controlled supervision-scaling experiment will test whether the relative
performance of shared and specialized adaptation changes with task-specific
data availability.

For selected tasks, training subsets will be constructed using

\begin{equation}
p\in\{1,5,10,25,50,100\}\%
\end{equation}

of the available training data.

At each data scale, we will compare:

\begin{enumerate}
    \item Shared LoRA,
    \item HydraLoRA,
    \item TaskLoRA.
\end{enumerate}

The primary experimental question is:

\begin{quote}
Does the preferred degree of parameter sharing change systematically with
the amount of task-specific supervision?
\end{quote}

Multiple random seeds will be used where computationally feasible.

The experiment will first be performed on a high-resource task through
controlled downsampling and subsequently evaluated on at least one naturally
lower-resource task.


% ============================================================
\section{Study 4: Partial and Soft Sharing}
\label{sec:study4}

\subsection{HydraLoRA}

HydraLoRA will be compared with Shared LoRA and TaskLoRA using matched
training data, optimization settings, parameter budgets where possible, and
evaluation splits.

The experiment will determine whether sharing only part of the low-rank
adaptation provides a different trade-off from the two endpoints.

\subsection{Soft-LoRA}

Weighted adapter composition will be evaluated using combinations of
task-specific adapters.

The analysis will examine whether the benefit of combining experts is
associated with the direct source--target transfer scores measured in
Study~\ref{sec:study2}.


% ============================================================
\section{Study 5: Modality and Domain Specialization}
\label{sec:study5}

\subsection{Modality-Level Specialization}

Modality-specific adapters will be trained independently.

We will compare:

\begin{enumerate}
    \item task-specific adapter only,
    \item modality-specific adapter only,
    \item task + modality adapter composition,
    \item learned task--modality composition.
\end{enumerate}

The objective is to test whether imaging modality provides a useful basis
for parameter specialization.

\subsection{Domain-Level Specialization}

Within classification, we will additionally compare:

\begin{enumerate}
    \item one global classification adapter,
    \item domain- or dataset-specific classification adapters,
    \item routed combinations of global and specialized adapters.
\end{enumerate}

The analysis will examine whether specialization below the task level is
useful when datasets share the same output structure but differ in medical
domain.


% ============================================================
\section{Study 6: Latent Expert Specialization}
\label{sec:study6}

We will evaluate latent specialization approaches including:

\begin{enumerate}
    \item LoRA mixture-of-experts,
    \item spectral expert adaptation.
\end{enumerate}

The following quantities will be measured:

\begin{itemize}
    \item expert utilization,
    \item routing entropy,
    \item task-conditioned expert probabilities,
    \item top-$k$ expert selections,
    \item pairwise distance between task-conditioned routing distributions,
    \item downstream task performance.
\end{itemize}

For tasks $t_i$ and $t_j$, routing-distribution distance can be defined as

\begin{equation}
D_{ij}
=
\frac{1}{2}
\sum_{k=1}^{K}
\left|
p(k\mid t_i)-p(k\mid t_j)
\right|.
\end{equation}

The study will test whether learned routing produces semantically distinct
task or capability specialization without explicit task assignment.


% ============================================================
\subsection{Router Perturbation Analysis}

Where applicable, the learned routing distribution $\pi$ will be gradually
perturbed toward a uniform distribution:

\begin{equation}
\tilde{\pi}
=
(1-\lambda)\pi+\lambda u,
\end{equation}

where $u$ is the uniform expert distribution and
$\lambda\in[0,1]$ controls perturbation strength.

This experiment will test how strongly downstream predictions depend on the
learned routing decisions.


% ============================================================
\section{Study 7: Representation Analysis}
\label{sec:study7}

To investigate the internal effect of different adaptation strategies,
hidden representations will be extracted from selected layers of:

\begin{itemize}
    \item the base VLM,
    \item Shared LoRA,
    \item TaskLoRA,
    \item HydraLoRA,
    \item Soft-LoRA,
    \item LoRA-MoE,
    \item spectral expert models.
\end{itemize}

Centered Kernel Alignment (CKA) and/or Representational Similarity Analysis
(RSA) will be used to compare learned representations.

For representation matrices $X$ and $Y$, CKA can be expressed as

\begin{equation}
\mathrm{CKA}(X,Y)
=
\frac{
\mathrm{HSIC}(X,Y)
}{
\sqrt{
\mathrm{HSIC}(X,X)
\mathrm{HSIC}(Y,Y)
}
}.
\end{equation}

We will examine whether representation similarity is associated with
cross-task transfer:

\begin{equation}
\rho
=
\mathrm{corr}
\left(
S_{ij},
C_{ij}
\right),
\end{equation}

where $S_{ij}$ denotes representation similarity between tasks $i$ and $j$
and $C_{ij}$ denotes the corresponding transfer score.

The analysis will address:

\begin{itemize}
    \item whether transferable tasks have more similar representations,
    \item whether shared adaptation preserves greater cross-task similarity,
    \item whether task specialization emerges gradually across layers,
    \item whether latent experts produce representationally distinct
    subspaces.
\end{itemize}


% ============================================================
\section{Study 8: Specialization Location}
\label{sec:study8}

Specialization will be compared at different locations within the VLM:

\begin{enumerate}
    \item vision encoder,
    \item multimodal projection pathway,
    \item late language-model blocks,
    \item LoRA adapters.
\end{enumerate}

Both downstream performance and representation shift will be evaluated.

The primary question is:

\begin{quote}
Which parts of a pretrained VLM can be specialized while preserving useful
pretrained multimodal representations?
\end{quote}


% ============================================================
\section{Study 9: Practical Adaptation Factors}
\label{sec:study9}

This section will investigate practical factors that may interact with the
adaptation strategy. Existing experiments will be included only when their
protocols permit scientifically meaningful comparison.

\subsection{Image Resolution}

Matched training and inference resolutions will be compared across selected
tasks and domains.

\subsection{LoRA Rank}

Task-specific rank sweeps will be used to investigate adapter-capacity
requirements.

\subsection{Class Imbalance}

Sampling, rebalancing, and related strategies will be evaluated for
classification and multi-label classification.

\subsection{Image Preprocessing}

Preprocessing experiments will include medically relevant intensity
transformations, normalization, cropping, and other domain-specific
operations where applicable.

\subsection{Data Augmentation}

Augmentation strategies will be evaluated according to whether they
represent plausible invariances for the corresponding medical imaging
domain.

\subsection{Multi-Image Handling}

Single-image and multi-view input strategies will be compared for datasets
containing multiple views or related images.

\subsection{Prompting}

Task-specific prompt formulations will be evaluated under controlled
settings.

\subsection{Decoding}

Greedy and stochastic decoding strategies will be compared for tasks where
generation settings may influence quantitative or structured outputs.


% ============================================================
\section{Study 10: Cross-Backbone Validation}
\label{sec:study10}

The core experimental questions will be evaluated using one secondary VLM.

The secondary-backbone experiments will include:

\begin{enumerate}
    \item Shared LoRA versus TaskLoRA,
    \item TaskLoRA versus HydraLoRA,
    \item a reduced supervision-scaling experiment.
\end{enumerate}

If the secondary backbone is medically pretrained, an additional analysis
will examine:

\begin{quote}
Does medical-domain pretraining change the relative need for downstream
task specialization?
\end{quote}

For each backbone $b$, the relative specialization effect can be summarized
as

\begin{equation}
\Delta_{\mathrm{spec}}^{(b)}
=
M_{\mathrm{TaskLoRA}}^{(b)}
-
M_{\mathrm{SharedLoRA}}^{(b)}.
\end{equation}

The comparison will be interpreted only after accounting for differences in
backbone architecture, parameter count, and training configuration.


% ============================================================
\section{Statistical Analysis}
\label{sec:statistics}

For principal experiments, we will report:

\begin{itemize}
    \item mean performance,
    \item 95\% bootstrap confidence intervals,
    \item paired bootstrap intervals for method differences,
    \item multiple random seeds for newly designed controlled experiments
    where feasible,
    \item correction for multiple comparisons where appropriate.
\end{itemize}

For the supervision-scaling study, the interaction between training-set size
and adaptation strategy will be tested.

A possible model is

\begin{equation}
Y
=
\beta_0
+
\beta_1\log N
+
\beta_2S
+
\beta_3(\log N\times S)
+
u_{\mathrm{task}}
+
\epsilon,
\end{equation}

where $N$ denotes the number of training samples and $S$ denotes the
adaptation strategy.

For representation analysis, statistical testing will assess the association
between representation similarity and cross-task transfer.


% ============================================================
\section{Planned Figures and Tables}
\label{sec:figures}

\subsection{Main Figures}

\begin{enumerate}

    \item \textbf{Sharing--specialization spectrum.}

    Conceptual organization of Shared LoRA, HydraLoRA, Soft-LoRA,
    latent experts, TaskLoRA, and domain specialists.

    \item \textbf{Experimental landscape.}

    Overview of datasets, modalities, task families, and output structures.

    \item \textbf{Cross-task transfer matrix.}

    Source adapters on rows and target tasks on columns.

    \item \textbf{Supervision-scaling curves.}

    Shared LoRA, HydraLoRA, and TaskLoRA across different amounts of
    task-specific supervision.

    \item \textbf{Latent expert routing analysis.}

    Task-conditioned expert-routing heatmaps and utilization statistics.

    \item \textbf{Representation similarity and transfer.}

    CKA/RSA similarity matrices together with their relationship to
    cross-task transfer.

    \item \textbf{Specialization granularity.}

    Comparison of global, task-, modality-, and domain-level specialization.

\end{enumerate}


% ============================================================
\subsection{Main Tables}

\begin{enumerate}

    \item Dataset, modality, task, sample-count, and evaluation-metric
    summary.

    \item Shared LoRA versus TaskLoRA versus HydraLoRA.

    \item Partial-sharing and latent-expert methods under matched
    experimental conditions.

    \item Supervision-scaling results.

    \item Task-, modality-, and domain-level specialization comparison.

    \item Cross-backbone validation.

\end{enumerate}


% ============================================================
\section{Synthesis of Experimental Evidence}
\label{sec:synthesis}

After completing the controlled studies, the experimental evidence will be
synthesized to examine whether consistent patterns emerge regarding:

\begin{itemize}

    \item when parameter sharing is beneficial,

    \item when explicit specialization is beneficial,

    \item whether supervision availability affects the relative value of
    sharing and specialization,

    \item whether output structure is associated with cross-task transfer,

    \item whether modality or medical domain predicts useful specialization,

    \item whether latent routing discovers meaningful expert structure,

    \item whether representation similarity predicts positive or negative
    transfer,

    \item where specialization can be introduced without adversely affecting
    pretrained multimodal representations,

    \item how practical model and inference configurations interact with
    specialization.

\end{itemize}

Any design recommendations will be formulated only from patterns that are
consistently supported by the controlled experiments.


% ============================================================
\section{Discussion}
\label{sec:discussion}

The discussion will interpret experimentally supported relationships among:

\begin{itemize}
    \item parameter sharing,
    \item specialization,
    \item task compatibility,
    \item data availability,
    \item output heterogeneity,
    \item imaging modality,
    \item medical domain,
    \item representation similarity.
\end{itemize}

The discussion will clearly distinguish empirical observations from possible
mechanistic explanations.

Potential explanations arising from the experiments will be treated as
hypotheses unless supported by additional analysis.


% ============================================================
\section{Limitations}
\label{sec:limitations}

The final paper will discuss limitations including:

\begin{itemize}

    \item dependence on the selected VLM backbones,

    \item variation in task-specific sample sizes,

    \item use of benchmark-derived rather than prospective clinical data,

    \item possible overlap between foundation-model pretraining data and
    evaluation datasets,

    \item imperfect comparability of historical experiments unless rerun
    under matched conditions,

    \item computational limitations on repeating every experiment across
    multiple VLM backbones,

    \item limitations of automated report-generation evaluation,

    \item coarse definitions of modality and medical domain.

\end{itemize}


% ============================================================
\section{Conclusion}
\label{sec:conclusion}

This study is designed to systematically investigate how parameter sharing
and specialization should be structured when adapting vision--language
models to heterogeneous medical image-analysis tasks.

By comparing multiple parameter-efficient adaptation strategies under
controlled conditions and analyzing cross-task transfer, supervision
scaling, routing behavior, representation similarity, specialization
granularity, specialization location, and practical model configuration,
the study aims to identify the factors associated with successful medical
VLM adaptation.

The final conclusions and any resulting design recommendations will be
restricted to experimentally verified results.


% ============================================================
\appendices

\section{Additional Experimental Details}

The appendix will contain:

\begin{itemize}
    \item complete training hyperparameters,
    \item full source--target transfer matrices,
    \item dataset-level results,
    \item complete routing statistics,
    \item additional representation analyses,
    \item complete supervision-scaling results,
    \item prompt variants,
    \item preprocessing configurations,
    \item augmentation experiments,
    \item rank sweeps,
    \item additional MoE configurations,
    \item additional backbone experiments.
\end{itemize}


% ============================================================
\section*{Acknowledgment}

Funding and acknowledgment information will be added before submission.


% ============================================================
\begin{thebibliography}{00}

% References will be added as the manuscript is developed.

% Examples:
%
% \bibitem{b1}
% E. J. Hu et al.,
% ``LoRA: Low-rank adaptation of large language models,''
% in Proc. Int. Conf. Learn. Representations (ICLR), 2022.
%
% \bibitem{b2}
% Y. Lai et al.,
% ``Med-R1: Reinforcement learning for generalizable medical reasoning
% in vision-language models,''
% IEEE Trans. Med. Imag., 2026.
%
% Add MedHallTune, Med-MoE, HydraLoRA, LoRAMoE,
% Spectral LoRA, generalist-vs-specialist VLM papers,
% and other related work here.

\end{thebibliography}

\end{document}
